\Page{15}{\textenglish{03 / CAMPAIGNS}}{מאיה לור צריכה להעלות עוד פוסט}{46}{52}

\Lead{למאיה לור יש מה להוציא מהתיק. היא יכולה לספר איזה מוצר טיפוח היא לוקחת איתה, לשחק עם שגרת בוקר או להגיב לדבר שהפיד התעכב עליו אתמול. הפרטים האלה הם חלק מהדמות שבנתה סופר־פארם, שהשיקה אותה בנובמבר 2025. לפי דיווח ההשקה ב־\textenglish{mako}, מומחי קריאייטיב, עיצוב, אנימציה ובינה מלאכותית עבדו על הדמות ועל הבעות, תנועה ודיבור. ניהול העמוד נועד להישאר בידי מחלקת הפרסום של הרשת. המקורות אינם מזהים את המודלים המסוימים ששימשו בפועל. \SourceRef{10}}

{\fontsize{14.2}{21.3}\selectfont

מנהלת התקשורת השיווקית רונה ליכטן־אוחיון סיפרה ל״גלובס״ על ישיבות אפיון שקבעו איך הדמות תיראה, תדבר ותופיע לאורך זמן. הרעיון השיווקי היה לפנות לקהל צעיר בשפה שמוכרת לו מן הפיד. תחילה הדמות נועדה לקדם את הרשת, ובהמשך נבחנו שיתופי פעולה. איגוד השיווק הישראלי תיאר פעילות גם אחרי ההשקה ומנה את המהלך בין זוכי פרס החדשנות לרבעון הרביעי של 2025. ההכרה מעידה על האופן שבו התקבל המהלך בוועדה מקצועית. \SourceRef{10}

\par}

\Panel{״גלובס״ דיווח על השקעה של שני מיליון שקלים ביצירת המשפיענית. הדיווח לא כלל פירוט סעיפים מבוקר, ולכן אפשר להיעזר בו כדי להבין את סדר הגודל שדווח למהלך, בלי להעמידו מול תשלום לדוגמנית או עלות של יום צילום. הדמות זקוקה לקול עקבי, עריכה ובקרה לאורך זמן. כל פוסט מכניס שוב לתמונה אנשים שמחליטים מה מתאים לדמות ולמותג. \SourceRef{10}}

{\fontsize{13.8}{21.3}\selectfont

האתגר הוא לשמור על העניין בלי לטשטש מי עומדת מאחורי ההמלצה. בטור דעה ב־\textenglish{mako} כתבה נטע חוטר שהליטוש והשליטה של המותג מקשים עליה להזדהות עם מאיה. ליכטן־אוחיון, מנגד, סיפרה שקיבלה תגובות אוהדות לצד שאלות על הטעם לעקוב אחר דמות שאינה קיימת. הטור והריאיון מתעדים תגובות שונות, בלי למדוד מי מהן מייצגת את העוקבות. מאיה יכולה להציג לוק חדש מחר בבוקר. הסיבה שבגללה הקהל יחזור לראות אותו תלויה בסיפור שהצוות בונה לאורך זמן. \SourceRef{10}

\par}
